\documentclass[11pt]{article}
\usepackage[margin=1in]{geometry}
\usepackage[T1]{fontenc}
\usepackage{amsmath,array,float,longtable}
\newfloat{algorithm}{tbp}{loa}
\floatname{algorithm}{Algorithm}
\renewcommand{\topfraction}{0.9}
\renewcommand{\bottomfraction}{0.8}
\renewcommand{\textfraction}{0.07}
\renewcommand{\floatpagefraction}{0.8}
\usepackage[hidelinks]{hyperref}
\setlength{\parindent}{0pt}
\setlength{\parskip}{5pt}
\newcommand{\code}[1]{\texttt{#1}}
\newcommand{\flowbox}[2]{\fbox{\parbox{0.79\linewidth}{\centering\strut\textbf{#1}\\#2\strut}}}
\title{Concord: from changed sources to maintained findings}
\author{}
\date{Proposed design}
\begin{document}
\maketitle

Akmon extracts a shared graph from source material. Concord maintains consistency
findings over that graph. A source update retracts old claims, inserts replacements
and re-checks the affected claim pairs. Optional repair selects claims to accept; an
approved source edit is a separate action.

The three finite algorithms below specify a bucket-local pair-rule core. They
make state and failure handling explicit. They refine the
Concord design (internal pull request 2521); they do not claim
that the general query engine or this complete integration is implemented.

Incremental maintenance is the core of the design. Algorithms 1 and 2 read only
the changed sources, the claims they remove or add, and the newly arrived model
answers, and they revisit only the claim pairs these can affect. Their work does not
grow with the rest of the graph. Three changes are broad by nature: a new
producer configuration changes every extraction request, a new rule revision
changes every finding key, and a change to a shared dictionary or policy
re-extracts every source that reads it. Answer cleanup, which keeps the answer
of every active source's current request, is a separate background pass over
the stored answers.

Section~\ref{sec:status} separates three things for every claim this document
makes: \emph{design intent} (what the design requires), \emph{proved model}
(what a checked Lean proof establishes about a simplified model), and
\emph{implementation agreement} (whether code matches the model). No Concord
code yet implements these algorithms.

\section{Algorithmic design}

\subsection{Architecture}

\begin{figure}[h]
\centering
\begin{tabular}{c}
\flowbox{Source items}{IDs, revisions, content and context}\\[3pt]
$\downarrow$\quad old/new source inputs\\[3pt]
\flowbox{Existing AKMON: shared extraction}{Algorithm 1 proposes finer per-item replacement}\\[3pt]
$\downarrow$\quad derived facts applied; publisher commits\\[3pt]
\flowbox{Pinned graph snapshots}{Claims and provenance; explicit coverage is required}\\[3pt]
$\downarrow$\quad signed claim changes and coverage\\[3pt]
\flowbox{Proposed Concord: maintenance}{Algorithm 2 re-tests affected claim pairs}\\[3pt]
$\downarrow$\quad findings, statuses and witnesses\\[3pt]
\flowbox{Stored Repair proposals}{Algorithm 3 persists acceptance choices; application is separate}\\
\end{tabular}
\caption{Shared extraction publishes claims once for all rules. Fixed stored answers
feed extraction and judged rules; extraction assigns each claim its bucket, and rules supply the tests.
Repair also needs the same pinned claims and protection policy. Human-approved
source correction returns to the first box. Applying an acceptance proposal never edits a source; source-edit proposals require an authorized edit.}
\end{figure}

Coverage travels with claims. A zero-row relation alone cannot establish that
extraction completed. Pinned answers are inputs to a run; an independent worker
creates observations for a later run. The diagram distinguishes the existing
Akmon substrate from proposed per-item coverage and Concord maintenance.


\subsection{Five graph records}

A \textbf{Source} is an identified paragraph, table or code declaration. It stores
source\_id, revision, content and context. For each producer at that revision it
also records extraction completion: complete, pending or failed, with receipt
links. Complete extraction may produce zero claims. Pending or failed extraction
cannot justify an all-clear.

A \textbf{Claim} stores id, bucket, value and provenance. It links to its Source
revision, supporting span and extraction receipt. Its ID includes the complete
extraction request and local output ID. This identity is safe because a
request's stored answer is written once and never replaced; a different answer
needs a different request, for example a new retry policy or producer
revision. Answer retention therefore keeps the answer of every active source's
current request. A bucket is (entity, predicate, scope),
or Unknown. Different source supports remain different claims; extraction does
not establish their truth.

A \textbf{Rule} stores id, revision, fixed context and a versioned test. The
pair-rule core below compares unordered claim pairs within a bucket. Outcomes
are NoViolation, Violation(conclusion), Pending(reason) or Failed(reason).
Unknown eligibility is pending, not false.

A \textbf{Finding} links its Rule revision, witness Claims and their Source
provenance, and stores the conclusion and evaluation identity. Its key contains
rule ID/revision and the ordered witness claim IDs; its stable name, to which
reviewer decisions attach, is defined below. Check completion is recorded per
rule and pair, so a bucket with no finding is known to be checked. Rule-wide incomplete
coverage marks the rule ``may be out of date'' or ``could not check'' (defined
in Implementation); a completed finding need not mean the
sources agree.

\textbf{Finding names.} A claim ID contains its full request, so editing any
part of a witness's paragraph, or changing the producer, gives the claim a new
ID and the pair a new key. Reviewer decisions therefore attach to a separate,
stable \emph{finding name}: the rule ID, the bucket, and the unordered pair of
the witnesses' (source ID, value) (operator decision, 2026-10-01). An edit to a
witness's paragraph that keeps its value keeps the name and the decision. A
changed value is a different conflict and gets a new name. A decision records
the rule revision and each witness's source revision and span that the reviewer
saw; when either differs, the finding shows ``changed since decision'', as in
the Concord design (internal pull request 2521). Two
supports from the same source with the same value give the same name and share
one decision; the decision records the set of (source revision, span) it saw,
so either finding shows ``changed since decision'' only when its own spans are
not in that set. Values are compared by their canonical typed serialization.
When a value changes, the finding's name changes: the old decision is retired
with its name and stays visible in that name's history, and the new finding
starts undecided. Computing a name reads only the pair's two claims.

These names depart from the Concord design on two points, by the operator's
decision. The Concord design names a finding by rule, subject and witness and
never by value; here the value is part of the name, so a changed value is a new
conflict. The Concord design retires decisions at a new rule version; here the
name holds the rule ID, not its revision, so a decision carries across a rule
revision and shows ``changed since decision''.

A \textbf{Repair} is a separately stored proposal, not a property of a Finding.
It stores id; its pins (the bucket, the conflict component's claim IDs and
edges, the complete hard-rule identities, and the selection-policy identity
with the protected IDs in the component); action set; accepted IDs; supporting
Finding links; lifecycle status; and approval and application receipts.
This algorithm creates one proposal per acceptance choice: actions are
RejectClaim(id), interpreted in a separate acceptance view. Source and Claim
evidence remains stored. The record can also represent an authorized source-edit
proposal, but the acceptance enumerator does not generate or execute edits.

Repair status is Proposed, Approved, Applied, Rejected or Superseded.
Approval names a specific proposal and authorizes its actions; it does not
apply them. Application re-derives the pinned component at the current state,
compares it and the rule and policy pins, and applies the actions in the same
transaction. A mismatch blocks application and marks the proposal Superseded;
superseding happens at application or at the next enumeration for that bucket,
not during maintenance. Application records its actual
effects and receipt. Applied records remain historical when sources change.

\subsection{One complete example}

\textbf{Initial state.} S1 revision 1 and S2 revision 1 both state that Orion
launches June 10. Their claims claim\_s1r1 and claim\_s2r1 retain different
spans but share the bucket (Orion, launch-date, public). Rule R1 is
deterministic and forbids different dates in a bucket; it is also the hard rule
for repair. Both Sources record completed extraction. The bucket's stored output
has no finding, no issue and no request.

\textbf{Source update, answer pending.} S2 revision 2 states June 12. The graph
difference lists S2, so EXTRACT selects only S2. Its new request,
request\_s2r2, has no answer yet: EXTRACT removes claim\_s2r1, stores a Pending
record, adds S2 to \code{unfinished} and records request\_s2r2 in
\code{pending\_index}. UPDATE receives one removed claim and deletes its one
pair, with claim\_s1r1; the bucket now holds claim\_s1r1 alone. R1 is marked
``may be out of date'' because S2's answer is Missing; zero known conflicts is
not an all-clear. No other bucket or source is read.

\textbf{Answer arrives.} A worker stores the answer to request\_s2r2. In the
next run \code{arrived\_requests} contains it, \code{pending\_index} maps it to
S2, and EXTRACT re-selects S2 alone, decodes claim\_s2r2 and clears S2 from
\code{unfinished}. UPDATE receives one added claim and runs one pair test,
claim\_s1r1 against claim\_s2r2, which stores a finding with both spans. Its
name is (R1, the bucket, \{(S1, June 10), (S2, June 12)\}). UPDATE does not
decide which is true. R1 is deterministic, so its pair read no request and no
answer arrival could select it.

\textbf{Repair choices.} The bucket's conflict graph has one edge, so one
component \{claim\_s1r1, claim\_s2r2\}. With no protected claims, Minimal
returns two proposals: reject \code{\{claim\_s1r1\}} or reject
\code{\{claim\_s2r2\}}. Protecting claim\_s1r1 permits only rejecting
\code{\{claim\_s2r2\}}; protecting both makes repair Impossible. Each choice
becomes a stored Proposed Repair pinned to the component, R1 and the policy. A
person approves the proposal rejecting claim\_s2r2. A separate application
re-derives the component, finds it unchanged, marks the proposal Applied,
updates the acceptance view, and supersedes the sibling proposal. S2's June 12 content, its Claim and the Finding
remain evidence.

\textbf{Source correction.} An approved edit creates S2 revision 3 stating June
10. When its answer arrives, EXTRACT retracts claim\_s2r2 and inserts
claim\_s2r3. UPDATE deletes claim\_s2r2's pair and tests claim\_s2r3 against
claim\_s1r1; the pair agrees, so the finding and the component disappear. Had
S2 revision 3 changed only its wording and kept June 12, the new pair would
have the same finding name, and a reviewer's decision on it would carry over,
marked ``changed since decision'' because the source revision differs. The other proposal (reject
claim\_s1r1) was already superseded when the first one was applied; had none
been applied, the retirement step after UPDATE would supersede both, because
their pinned component no longer exists. The Applied proposal remains a
historical record. The bucket again has no finding and no issue.

\subsection{Algorithms over the stored graph}

The procedures below use finite maps as the storage representation of these
five records. The extraction map is a private index of Source completion and
source-supported Claims, not another graph object. Its request-cache fields,
decoder contracts and helper notation are specified in Implementation.

Each algorithm starts from stored state and a change, and returns new state.
Working maps are \emph{overlays}: \code{overlay(map)} reads through to the
stored map and records only the entries written, so starting one costs nothing
and its cost is the number of writes. All maps and indexes are committed
atomically with the new identity.

\subsection{Algorithm 1: maintain shared extraction records}

Algorithm~\ref{alg:extract} gives the procedure. \textit{Input:} the old
extraction map; its pending-request index; the set \code{unfinished}, which
holds the sources whose extraction is incomplete; the new
source-item map; the producer; the fixed answer snapshot; \code{changed\_sources},
the source IDs that this publication added, removed or edited; and
\code{arrived\_requests}, the requests whose answer is new in this snapshot.
\textit{Output:} the new extraction map and index, the claims removed and
added, and the new \code{unfinished} set.

\textit{Preconditions.} \code{changed\_sources} contains every source whose item
differs between the old and new source maps; it comes from the graph store's
difference between the two checkpoints. \code{pending\_index} maps each request
whose answer is Missing to the sources waiting on it. A request's answer only
moves from Missing to a stored value; it is never replaced (first answer wins).
Each stored record's claims and issues equal the decoding of its stored answer,
and the decoder and \code{make\_claim} read only fields covered by the request.
The extraction map stores the producer identity it was built with; a new
producer configuration changes every request, so lines 6--7 select every source
themselves rather than trusting the caller.

Two other inputs of a request are shared across sources, and both stay local.
A source's context records the revision of every dictionary or policy it reads,
and a stored index maps each such input to the sources that read it; when a
shared input changes, the graph difference lists that input and the index
expands it into \code{changed\_sources}. That pass is proportional to the
input's readers, which can be the whole graph for a dictionary every source
reads. Retrying a stored failure is an edit of that one source: the source's
item carries its attempt number, so a retry changes only its request.

\begin{algorithm}[tbp]
{\footnotesize
\begin{verbatim}
01 EXTRACT(old_extractions, pending_index, unfinished, new_sources,
02         producer, answers, changed_sources, arrived_requests):
03   selected = set(changed_sources)
04   for request in arrived_requests:
05     selected |= pending_index.get(request, set())
06   if producer.identity != old_extractions.producer:
07     selected = keys(old_extractions) | keys(new_sources)
08   extractions = overlay(old_extractions)
09   pending, unfinished = overlay(pending_index), overlay(unfinished)
10   removed, added = set(), set()
11   extractions.producer = producer.identity
12   for source_id in selected:
13     if source_id in old_extractions:
14       old = old_extractions[source_id]
15       removed |= old.claims
16       if old.answer is Missing:
17         pending[old.request].remove(source_id)
18       unfinished.discard(source_id)
19       extractions.pop(source_id, None)
20     if source_id not in new_sources: continue
21     request = make_request(source_id, new_sources[source_id],
22                            producer)
23     answer = answers.get(request, Missing)
24     claims, issues = set(), set()
25     if answer is Missing:
26       issues = {Pending(request)}
27       pending[request].add(source_id)
28     elif answer is Failure(reason): issues = {Failed(reason)}
29     else:
30       decoded = producer.decode(answer.payload)
31       if decoded is Invalid(reason): issues = {Failed(reason)}
32       else: claims = {make_claim(request, row,
33                         new_sources[source_id], answer)
34                       for row in decoded.rows}
35     if issues != set(): unfinished.add(source_id)
36     extractions[source_id] =
37       ExtractionRecord(request, answer, claims, issues)
38     added |= claims
39   return (extractions, pending, unfinished,
40           removed - added, added - removed)
\end{verbatim}}
\caption{EXTRACT maintains the shared extraction records.}\label{alg:extract}
\end{algorithm}

Replacement removes every prior claim and status of a selected source. Each
support has its own claim ID, so the removed and added sets retain signed
support changes. A claim ID contains its request, and a request contains its
source ID, so claims of different sources never coincide; the claims of
unselected sources are unchanged. Hence the returned sets equal the differences
between all old and all new claims, and \code{extractions} equals a rebuild from
scratch, while the work is proportional to the selected sources, their claims
and the arrived requests. A source that is re-selected without a real change
returns the same claims, which cancel in lines 39--40. A temporary worker
failure leaves the answer Missing, so the source stays pending until an outcome
is stored. A source also stays in \code{unfinished} when its answer was a stored
failure; Pending and Failed sources are both incomplete coverage.

\subsection{Algorithm 2: maintain inconsistency findings}

A pair's outcome depends only on its two claims, the rule, and the answers to
the requests its test reads. Algorithm 2 therefore maintains outcomes pair by
pair: it deletes the pairs of removed claims, tests the pairs of added claims,
and re-tests the pairs whose requests received answers. No other pair is
touched, so one changed claim in a bucket of $n$ claims costs at most $n$ pair
tests (operator decision, 2026-10-01).

Algorithm~\ref{alg:update} gives the procedure. \textit{Input:} the stored
check state; the claims removed and added by Algorithm 1;
\code{arrived\_requests}; the new rule and answer snapshot; and the new
identity. \textit{Output:} the check state for the new identity.

The check state holds, for one rule:
\begin{itemize}
\setlength{\itemsep}{0pt}
\item \code{claims\_by\_bucket}, each known bucket's member claims, and
  \code{unplaced}, the IDs of claims whose bucket is Unknown;
\item \code{pairs}, keyed by the ordered pair of claim IDs (smaller first), with
  each pair's bucket, its two claims, its outcome and the requests its test
  read;
\item \code{pairs\_by\_request}, the reverse index from each request to the
  pairs whose test read it;
\item \code{violations\_by\_bucket} and \code{issues\_by\_bucket}, the keys of
  each bucket's pairs whose outcome is a Violation, or Pending or Failed.
\end{itemize}
The \emph{full reference check} tests every pair of every known bucket with the
same \code{check\_pair}. The helper \code{pair\_key(a, b)} orders two claim
IDs. The helper \code{store(key, pair)} writes a pair and adds it to the request index
and to the bucket index for its outcome; \code{drop(key)} removes a pair and
all its index entries. Both cost the pair's request count. The two bucket
indexes are overlays too, started like \code{pairs\_by\_request} on line 6.
UPDATE also returns \code{touched}, the buckets of the removed and added known
claims and of the re-tested pairs (lines 37--38); these are the \emph{affected
buckets} that repair retirement reads. A rule change makes every bucket
affected.

\textit{Precondition.} The stored state equals the full reference check of the
old claims, rule and answers; every removed claim is an old claim; no added claim
is already among the old claims that remain; and claim IDs are unique (Algorithm 1's (request, local ID) IDs are, and its removed set
is drawn from the old claims). \code{arrived\_requests} contains every request
whose answer differs between the old and new snapshots; under first answer
wins, these are the answers stored since the old snapshot. UPDATE refuses, as
InvalidInput, a stored state whose identity is not the base checkpoint, rule
identity and answer sequence that the caller's change starts from (line 3).

\begin{algorithm}[tbp]
{\footnotesize
\begin{verbatim}
01 UPDATE(state, base_identity, removed_claims, added_claims,
02        arrived_requests, new_rule, new_answers, new_identity):
03   if state.identity != base_identity: return InvalidInput
04   claims_by_bucket = overlay(state.claims_by_bucket)
05   unplaced, pairs = overlay(state.unplaced), overlay(state.pairs)
06   pairs_by_request = overlay(state.pairs_by_request)
07   for claim in removed_claims:
08     if claim.bucket == Unknown: unplaced.discard(claim.id)
09     else:
10       claims_by_bucket[claim.bucket].discard(claim)
11       for other in claims_by_bucket[claim.bucket]:
12         drop(pair_key(claim.id, other.id))
13   if new_rule.identity != state.rule_identity:
14     return FULL_CHECK(claims_by_bucket, unplaced, added_claims,
15                       new_rule, new_answers, new_identity)
16   tested = set()
17   for claim in added_claims:
18     if claim.bucket == Unknown: unplaced.add(claim.id); continue
19     for other in claims_by_bucket[claim.bucket]:
20       left, right = ordered(claim, other)
21       outcome, requests = check_pair(new_rule, left, right,
22                                      new_answers)
23       store(pair_key(left.id, right.id), Pair(claim.bucket, left,
24             right, outcome, requests))
25       tested.add(pair_key(left.id, right.id))
26     claims_by_bucket[claim.bucket].add(claim)
27   for request in arrived_requests:
28     for key in state.pairs_by_request.get(request, set()):
29       if key in pairs and key not in tested:
30         old = pairs[key]
31         outcome, requests = check_pair(new_rule, old.left,
32                                        old.right, new_answers)
33         drop(key)
34         store(key, Pair(old.bucket, old.left, old.right,
35                         outcome, requests))
36         tested.add(key)
37   touched = {c.bucket for c in removed_claims | added_claims
38              if c.bucket != Unknown} | {pairs[k].bucket for k in tested}
39   return CheckState(new_identity, new_rule.identity,
40     claims_by_bucket, unplaced, pairs, pairs_by_request,
41     violations_by_bucket, issues_by_bucket), touched
\end{verbatim}}
\caption{UPDATE maintains one rule's pair outcomes.}\label{alg:update}
\end{algorithm}

\textbf{Which pairs are touched.} Removing a claim deletes its pairs with every
claim still in its bucket (lines 7--12); when two removed claims share a
bucket, the first removal deletes the pair between them. Adding a claim tests
it against every claim already in its bucket, including claims added earlier in
the same run (lines 17--26), so each new pair is tested exactly once. A claim
that moves from bucket $b$ to $b'$ is a removed claim in $b$ and an added claim
in $b'$. A claim entering or leaving Unknown changes only \code{unplaced} and
the pairs of the one known bucket it leaves or enters. An arrived answer
re-tests exactly the surviving pairs whose test read its request (lines
27--36); a Pending outcome therefore clears when, and only when, its answer
arrives. Nothing else is revisited. A pending extraction elsewhere, or an
unplaced claim, is reported as rule-wide incomplete coverage, which blocks an
all-clear but changes no pair.

\textbf{Why the result is exact.} Every pair of the new claims either consists
of two old claims or contains an added claim. A pair of two old claims that
are both still present was not deleted; if none of its requests arrived, its
claims and the answers it reads are unchanged, so its stored outcome is still
the full reference outcome, and if one arrived, lines 27--36 re-test it. Every
pair containing an added claim is tested in lines 17--26. Every stored pair
containing a removed claim was deleted in lines 7--12. A rule change re-tests
every pair (\code{FULL\_CHECK}: drop all pairs, then test every pair of every
bucket after applying the claim changes); this pass is inherent, since every
finding key contains the rule revision. Section~\ref{sec:status} names the
Lean theorems that prove this for a model of the algorithm.

\textbf{Cost.} Outside a rule change, the work is
\[
\sum_{x\in\text{removed}} \bigl(n_{b(x)} + d_x\bigr)
\;+\; \sum_{y\in\text{added}} \bigl(n_{b(y)} + r_y\bigr)
\;+\; \sum_{q\in\text{arrived}} \bigl(|\text{pairs\_by\_request}[q]| + r_q\bigr)
\;+\; 1,
\]
where $n_{b}$ is the number of claims in bucket $b$, and $d$, $r$ count the
request-index entries deleted and written. One changed claim costs at most
$n_b$ pair tests, not $\binom{n_b}{2}$. This keeps within reach the work
gate of the Concord design (internal pull request 2521), one eighth of a
full recompute, even when one bucket holds most of
a rule's claims. No term depends on pairs that contain no changed claim and
read no arrived request.

\subsection{Algorithm 3: propose acceptance repairs}

Repair works on the findings that Algorithm 2 maintains, not on a fresh search
of the graph. In this version its hard constraints are pair rules only
(operator decision, 2026-10-01). Constraints over three or more claims, and
global contradictions, are not supported; a constraint that spans buckets would
also join components across buckets and lose the locality below. Within one bucket, the
Violation findings of the hard rules form a \emph{conflict graph}: claims are
vertices and each violating pair is an edge. A set of rejected claims is
\emph{safe} if it contains at least one end of every edge, so that no violating
pair remains among the accepted claims. Edges never cross buckets, so the
conflict graph splits into connected components inside buckets, and a
component's choices do not interact with any other component.

Algorithm~\ref{alg:repair} gives the procedure. \textit{Input:} a bucket and
one connected component of its conflict graph (its claim IDs and edges), the
hard rules' check states, the current identity (graph checkpoint and answer
sequence), the protected claim IDs, the selection policy, and whether the hard
rules' coverage is incomplete. \textit{Output:} a RepairEnumeration.

A policy contains revision, mode and weights. Modes are Minimal or Weighted.
Weighted mode requires a finite positive rational weight for every claim in the
component; \code{valid\_weight(value)} tests precisely that condition. Positive
weights matter: a zero weight would let a minimum-cost choice reject a claim for
no reason. A RepairEnumeration contains status, proposals and reason;
non-Available results contain no proposals, and
\code{repair\_result(status, proposals, reason)} constructs it.
\code{connected\_components(edges)} returns each component as its claim IDs and
edges. The caller passes a component that \code{COMPONENTS} returned for the
current state; application re-derives it, so a stale component cannot be
applied. \code{weighted} abbreviates \code{policy.mode == Weighted}; Minimal
proposals pin no weights, so a weight change does not supersede them. \code{ready(bucket, hard\_states, current)} is lines 9--13: it holds
when every hard rule's state is at the current identity and none of their
pairs in the bucket has a Pending or Failed outcome. \code{make\_repair}
constructs a Proposed Repair whose actions are RejectClaim(id) for each
rejected ID and whose ID is the canonical serialization of its pins and
actions; \code{persist\_if\_absent} stores it by ID, adds it to
\code{proposals\_by\_bucket}, and never overwrites existing lifecycle
receipts, so a repeated run returns the existing proposal.

\begin{algorithm}[tbp]
{\footnotesize
\begin{verbatim}
01 COMPONENTS(bucket, hard_states):
02   edges = set()
03   for state in hard_states:
04     for key in state.violations_by_bucket.get(bucket, set()):
05       edges.add(frozenset(key))
06   return connected_components(edges)

07 REPAIR(bucket, component, edges, hard_states, current,
08        protected_ids, policy, coverage_incomplete):
09   for state in hard_states:
10     if (state.identity.checkpoint, state.identity.answers) != current:
11       return repair_result(Unavailable, set(), "stale check")
12     if state.issues_by_bucket.get(bucket, set()) != set():
13       return repair_result(Unavailable, set(), "unfinished check")
14   if policy.mode not in {Minimal, Weighted}:
15     return repair_result(InvalidPolicy, set(), "invalid mode")
16   if policy.mode == Weighted and any(
17       not valid_weight(policy.weights.get(id)) for id in component):
18     return repair_result(InvalidPolicy, set(), "invalid weights")
19   safe = set()
20   for rejected in all_subsets(component - protected_ids):
21     if all(edge & rejected != set() for edge in edges):
22       safe.add(rejected)
23   if safe == set():
24     return repair_result(Impossible, set(), "protected conflict")
25   if policy.mode == Minimal:
26     selected = {rejected for rejected in safe
27                 if not any(other < rejected for other in safe)}
28   else:
29     costs = {rejected: sum(policy.weights[id] for id in rejected)
30              for rejected in safe}
31     selected = {rejected for rejected in safe
32                 if costs[rejected] == min(costs.values())}
33   pins = Pins(bucket, component, edges,
34     {state.rule_identity for state in hard_states},
35     (policy.revision, policy.mode,
36      {id: policy.weights[id] for id in component} if weighted else {}),
37     protected_ids & component)
38   proposals = set()
39   for rejected in selected:
40     proposals.add(persist_if_absent(make_repair(pins, rejected)))
41   return repair_result(Available, proposals,
42     "incomplete coverage" if coverage_incomplete else "")
\end{verbatim}}
\caption{COMPONENTS and REPAIR propose acceptance repairs for one conflict component.}\label{alg:repair}
\end{algorithm}

\textbf{What it guarantees.} Every proposal is safe and never rejects a
protected claim (lines 19--22). In Minimal mode the proposals are exactly the
inclusion-minimal safe sets of the component; in Weighted mode they are exactly
the minimum-weight safe sets. Because components do not interact, a minimal (or
minimum-weight) safe set for the whole bucket is exactly one such choice per
component, so per-component proposals represent every whole-bucket choice
without materializing the product. Repair runs only on current, finished
checks (lines 9--13): a hard rule whose update has not yet reached the current
checkpoint could be missing an edge, and a Pending or Failed pair could be an
unknown conflict, so either makes the bucket's repair Unavailable. An unknown
conflict is never treated as no conflict. Repair does not choose a true source,
and evidence is preserved. Certainty requires agreement across all selected
acceptance sets and, separately, admitted interpretations.

\textbf{Cost.} Lines 1--6 read one bucket's violation index, and lines 9--13
read one identity and one issue index per hard rule. Lines 19--22 enumerate $2^{k}$ subsets for a
component of $k$ unprotected claims; the search is exponential in the
component, not the graph, and the number of proposals can itself be
exponential in $k$. The pins hold only the component's weights and protected
IDs, so a weight change elsewhere does not touch a proposal. Incomplete coverage
is reported with the result, not stored in the proposal, because it changes
over time without changing the proposal: an unplaced claim, or a source still
being extracted, may later add a claim to the bucket, and application then
sees a different component.

\textbf{Application.} Application runs in one transaction. It requires
\code{ready(bucket, hard\_states, current)}, re-derives the component with
\code{COMPONENTS} for the pinned bucket, and requires the component, its edges,
the rule identities, the policy pin and the protected IDs to equal the pins.
If any check fails, the proposal becomes Superseded and the person is shown a
fresh enumeration. Otherwise it applies the actions, records their effects and
receipt, and marks every other Proposed or Approved repair pinned to the same
bucket and component Superseded, whatever its policy or protected IDs, because
the component now has its chosen acceptance. An Applied rejection stays in the
acceptance view until a person revokes it. If the component later changes, its
pins no longer match, so a fresh enumeration is offered for the new component;
it never silently reuses the old application. All of this
reads one bucket and its proposals. Unrelated updates elsewhere in the graph
do not invalidate a proposal.

\textbf{Retiring stale proposals.} After Algorithm 2 commits, the stored
proposals of each affected bucket, found through the index
\code{proposals\_by\_bucket}, are compared with the bucket's current
components; a Proposed or Approved repair whose pinned
component no longer exists becomes Superseded. This costs work proportional to
the affected buckets' proposals.

\section{Implementation}

Akmon already owns import, shared extraction, derived-fact application and graph
publication. Its ClaimExtractor reads pinned Paragraph/Theorem/Table records.
Current complete-producer replacement can cover broader output than Algorithm
1's proposed per-item state. Chronos supplies before/after checkpoint differences,
whose fallback may read entire snapshots. Current claim projection supports
limited bucket updates; these algorithms are a proposed precise pair-rule core.

ExtractionRecord is private implementation bookkeeping for shared extraction;
it is not a sixth conceptual graph object. Store the extraction map keyed by
source ID with its producer identity, the pending-request index, the
\code{unfinished} set, the index from each shared dictionary or policy to the
sources that read it, \code{proposals\_by\_bucket}, and, per
rule, the check state: \code{claims\_by\_bucket}, \code{unplaced},
\code{pairs}, \code{pairs\_by\_request}, \code{violations\_by\_bucket} and
\code{issues\_by\_bucket}, plus reviewer decisions keyed by finding name. These
maps represent Source completion and Rule pair completion; they are not additional
conceptual graph objects. Commit them, the Claims and the Findings with the
pinned source, rule and answer identity in one transaction. A journal resumes the
same target after interruption. Model fetching is outside evaluation; retain
answers required by committed stage inputs and pinned unfinished jobs,
including inputs that emitted no finding.

Algorithms 1 and 2 need three inputs from the substrate that are not free:
\code{changed\_sources} comes from Chronos's row difference between two
checkpoints (\code{diff\_rows}), which may fall back to reading whole snapshots;
\code{arrived\_requests} needs the answer store to list answers stored since a
snapshot, for example by insertion sequence; and the overlays need storage that
writes a delta rather than copying a map. If any of these is implemented by a
scan, that scan is part of the cost and the work is no longer local.

Totals, rankings, requirements and downstream stored relations require separate
query and dependency contracts; they are not silently executed by UPDATE's pair
loop. Repair is local to one conflict component only because its hard rules
are pair rules; a constraint that spans buckets would join components.

\subsection{Extraction and evaluation bookkeeping}

A source item is one independently identified paragraph, table or code
declaration. Its record has source\_id, revision, content and context. Context
includes every dictionary or policy input the producer reads. Producer
configuration contains code/decoder revision, model, prompt, schema and parameters.
\code{make\_request(source\_id, item, producer)} serializes only source\_id,
revision, content, context and producer configuration. It excludes derived
completion and receipt fields, so recording completion does not change its own
request.
An immutable answer map returns Missing, Success(payload) or Failure(reason);
a changed retry policy creates a different request.

The producer's total decoder validates field types, unique local IDs and source
spans. It returns Decoded(rows) or Invalid(reason). The Claim representation contains
id, bucket, value and provenance. Its ID is (complete request, local ID); its
bucket is (entity, predicate, scope), or Unknown if placement is unresolved.
\code{make\_claim(request, row, source, answer)} builds those fields: ID from
the request/local ID, bucket and value from the row, and provenance from the
source, span and answer receipt. Claim IDs are ordered by their
typed, length-prefixed serialization bytes. Extraction completion does not
establish factual truth.

An \textbf{ExtractionRecord} stores the extraction result for one source item:
\code{request}, \code{answer}, \code{claims}, and \code{issues}. It is not another
kind of graph object. The map \code{old\_extractions} is keyed by source\_id.
Completed extraction has an empty issue set, even when it produces zero claims.
For example, request\_s2r2 is the complete request for S2 revision 2, and
claim\_s2r2 is its June 12 claim. Its completed record is:

\begin{minipage}{\linewidth}{\small
\begin{verbatim}
ExtractionRecord(request=request_s2r2, answer=Success(payload),
                 claims={claim_s2r2}, issues=set())
\end{verbatim}}\end{minipage}

Before its answer is available, the same source item has:

\begin{minipage}{\linewidth}{\small
\begin{verbatim}
ExtractionRecord(request=request_s2r2, answer=Missing,
                 claims=set(), issues={Pending(request_s2r2)})
\end{verbatim}}\end{minipage}

The following names have one meaning throughout. \code{claim\_ids(claims)}
returns claim IDs, not claim records. Empty maps are \code{\{\}}; empty sets are
\code{set()}. In pseudocode, \code{|} means set union and \code{-} set
difference. \code{keys(map)} returns its key set. Rejection sets are immutable
mathematical values; pseudocode collections may contain them. \code{add(value)}
updates a working set by union with the singleton \code{\{value\}};
\code{remove(value)} and \code{discard(value)} delete it, and \code{discard}
is a no-op if it is absent. Indexing a map of sets at a missing key yields an
empty set, and a set left empty is deleted, so \code{keys} lists only non-empty
entries. \code{overlay(map)} is defined in Algorithms over the stored graph.
\code{pop(key, None)} removes an entry and is a no-op if absent. Assignment
uses \code{=}, equality uses \code{==}; \code{is Tag(fields)} matches a tagged
value and binds its fields; \code{\_} ignores a field. For sets, \code{<=}
means subset and \code{<} proper subset.

\emph{Rule-wide incomplete coverage} means that \code{unfinished} (Algorithm 1)
or \code{unplaced} (Algorithm 2) is non-empty. It blocks an all-clear but
changes no pair's outcome. Following the
Concord design (internal pull request 2521), it has two
kinds. A source whose answer is still Missing makes the rule's findings
\emph{may be out of date}; this clears when the answer arrives. A source whose
answer is a stored failure, or a claim whose completed extraction left its
bucket Unknown, makes the rule \emph{could not check}, naming those sources and
claims; this is permanent until the source changes. Keeping a count of each
kind makes both tests constant work.

The Rule representation contains id, revision, fixed context and a total pair test. It
compares every unordered pair in each known bucket. Test outcomes are
NoViolation, Violation(conclusion), Pending(reason) or Failed(reason).
Definitely ineligible pairs return NoViolation; unknown eligibility returns
Pending; errors return Failed.

\code{check\_pair(rule, left, right, answers)} returns (outcome, requests).
For deterministic rules, it calls the supplied test and returns an empty request
set. For judged rules, it serializes the complete rule, context and ordered claim
records into a request. Missing yields Pending; Failure yields Failed; Success
passes through the total typed decoder, with invalid output converted to Failed.
The request set contains that request even if absent. No answer is fetched here.
Complete rule identity includes immutable test/decoder code references.

A \textbf{CheckState} contains \code{identity}, \code{rule\_identity} and the
maps listed under Algorithm 2. A Pair has \code{bucket}, \code{left},
\code{right}, \code{outcome} and \code{requests}; \code{ordered(a, b)} returns
two claims with the smaller ID first. A bucket whose \code{issues\_by\_bucket}
entry is empty is completely evaluated, even with findings. The full FindingKey
is (rule\_id, rule\_revision, left\_id, right\_id): the rule part comes from
\code{rule\_identity}. A Finding is a pair whose outcome is a Violation; its
conclusion and its claims' provenance come from the stored pair, and its owning
CheckState supplies the Rule link and the evaluation identity. Reviewer
decisions attach to the finding name, not to this key. A retained older checkpoint is separate from this
current state.

The identity is (graph checkpoint, rule identity, answer sequence). The graph
checkpoint is the publication's checkpoint ID. The rule identity is the full
canonical typed serialization of the rule's ID, revision, context and code
references; it is not an assumed collision-free hash, and it is small. Because
answers are never replaced, the answers stored up to a sequence number form a
snapshot, so the sequence number identifies it. None of the three requires
reading the graph.


\section{Verification status}\label{sec:status}

Three levels are kept apart. \emph{Design intent} is what this document
requires. A \emph{proved model} is a simplified mathematical version of a
mechanism together with a Lean proof, checked by Lean 4.34.1 with no axioms
beyond \code{propext} and \code{Quot.sound}. \emph{Implementation agreement}
means production code is shown to behave like the model. No Concord code
implements Algorithms 1--3 yet, so no row below has implementation agreement.
Each proved model also has an executable Python control that enumerates small
cases and confirms that deliberately broken variants give wrong results.

\medskip
{\small
\begin{longtable}{@{}>{\raggedright}p{0.25\linewidth}>{\raggedright}p{0.37\linewidth}>{\raggedright\arraybackslash}p{0.30\linewidth}@{}}
\textbf{Design intent} & \textbf{Proved model} & \textbf{Not covered}\\\hline
\endhead
Algorithm 2 is exact: deleting removed claims' pairs, testing added claims'
pairs and re-testing pairs whose requests arrived gives the full check. &
\code{ConcordPairLevel.pair\_step\_exact} and \code{pair\_runs}: for any pair
rule that reads answers only through its reported requests, the whole stored
state (bucket members, unplaced IDs, every pair's outcome and requests, and a
complete and sound request index) stays exact across any number of runs;
\code{pair\_view\_exact}: each bucket's view equals the full check. Proved
counterexamples show stale results when a batch's added claims are not tested
against each other, when a removed claim's pairs are kept, or when arrived
answers are ignored. &
Rule changes (a full pass by design); calls whose answer chooses the next
request; real index and storage costs; the bucket indexes
\code{violations\_by\_bucket} and \code{issues\_by\_bucket} that repair
reads (the proved view rebuilds a bucket from its pairs). Assumes the removed
claims are old claims
and claim IDs are unique, which Algorithm 1's (request, local ID) IDs give but
no theorem links.\\[3pt]
Algorithm 2's work grows with the change, not the graph. &
\code{ConcordPairLevel.pair\_cost\_bound}: pair tests are at most the sum of
added claims' bucket sizes plus the index entries of arrived requests, and
deletions at most the sum of removed claims' bucket sizes;
\code{pair\_cost\_single}: one changed claim costs at most $n$ tests or
deletions in a bucket of $n$ claims. &
Counts tests and deletions, not index edits or measured time.\\[3pt]
Algorithm 1 equals a rebuild from scratch, and its removed and added sets are
exact. &
\code{ConcordExtract.extract\_exact}, \code{delta\_exact} and
\code{extract\_step}: from an exact stored state, one run yields the exact
state for the new sources and answers. &
Producer configuration changes (all sources selected); decoder determinism is
assumed.\\[3pt]
Algorithm 3 returns exactly the minimal or minimum-weight safe sets of a
component, and components can be repaired separately. &
\code{ConcordRepair.repair\_sound}, \code{repair\_complete} (both modes),
\code{component\_decomposition} (Minimal mode),
\code{min\_weight\_is\_minimal}. &
COMPONENTS, the \code{ready} test, pins, persistence, application and
retirement are not modeled. Weighted-mode decomposition is checked only on all
graphs of four claims. Output can be exponential.\\[3pt]
Algorithm 1's delta feeds Algorithm 2. &
Not proved: the two models use different claim types. &
The composition of Algorithms 1 and 2.\\[3pt]
Missing answers are reported, never treated as values. &
\code{unmarked\_sound} in \code{ConcordUncertainty}, for addition expressions. &
Unknown answers that decide a bucket, a ranking or a branch.\\[3pt]
Bounded lag under a work envelope. &
\code{ConcordBounded} FIFO backlog and deadline theorems. &
A real target's queue charge is assumed, not derived from Algorithm 2's
cost.\\[3pt]
Crash and resume keep progress without duplicate effects. &
\code{ConcordRecovery}, for an atomic journal and a deterministic evaluator. &
Database atomicity; provider faults.\\[3pt]
Answer cleanup keeps every answer current results need. &
\code{ConcordCleanup}, given complete roots. &
That the roots are complete for real stages.\\
\end{longtable}}

\medskip
The bucket-level model in \code{ConcordPairs} proves the same exactness for
the earlier whole-bucket variant and supplies the full check that
\code{ConcordPairLevel} is compared against. Two older proofs are weaker than
their names suggest. In \code{ConcordWork}
the set of touched groups is an input, so its locality theorems hold by the
definition of the cost; \code{ConcordPairs} replaces that assumption with a
derived set for the pair-rule core. In \code{ConcordProgress} the target's
queue charge is a hypothesis. The
W2 Lean work (internal pull request 2707) proves properties
of the broader Concord specification; the models here are separate and
smaller. Producer interpretation accuracy is a further, empirical question.
Every proof here is reproduced by
\begin{center}{\small\code{python3 formal/concord\_bounded\_maintenance/program/run\_checks.py}}\end{center}

\end{document}
